% ==========================================
% ABKÜRZUNGEN
% ==========================================

\newacronym{api}{API}{Application Programming Interface}
\newacronym{cicd}{CI/CD}{Continuous Integration/Continuous Deployment}
\newacronym{ide}{IDE}{Integrated Development Environment}
\newacronym{idp}{IdP}{Identity Provider}
\newacronym{jwt}{JWT}{JSON Web Token}
\newacronym{mvp}{MVP}{Minimum Viable Product}
\newacronym{oauth}{OAuth}{Open Authorization}
\newacronym{oidc}{OIDC}{OpenID Connect}
\newacronym{orm}{ORM}{Object-Relational Mapping}
\newacronym{pvc}{PVC}{Persistent Volume Claim}
\newacronym{dto}{DTO}{Data Transfer Object}
\newacronym{rest}{REST}{Representational State Transfer}
\newacronym{roi}{ROI}{Return on Investment}
\newacronym{sso}{SSO}{Single Sign-On}
\newacronym{uml}{UML}{Unified Modeling Language}

% ==========================================
% GLOSSAREINTRÄGE
% ==========================================

\newglossaryentry{backend}{
    name={Backend},
    description={Serverseitige Komponente einer Anwendung, die Geschäftslogik, Datenbankzugriffe und API-Endpunkte bereitstellt}
}

\newglossaryentry{clean_architecture}{
    name={Clean Architecture},
    description={Softwarearchitekturmuster, das eine klare Trennung von Geschäftslogik und technischen Details durch Schichtenbildung erreicht}
}

\newglossaryentry{devops}{
    name={DevOps},
    description={Kunstwort aus Development und Operations; beschreibt die Kombination von Softwareentwicklung und IT-Betrieb zur Verbesserung der Zusammenarbeit und Automatisierung}
}

\newglossaryentry{dsgvo}{
    name={DSGVO},
    description={Datenschutz-Grundverordnung der Europäischen Union (Verordnung (EU) 2016/679); regelt den Schutz personenbezogener Daten und deren Verarbeitung in der EU}
}

\newglossaryentry{frontend}{
    name={Frontend},
    description={Clientseitige Komponente einer Anwendung, die die Benutzeroberfläche und Benutzerinteraktion bereitstellt}
}

\newglossaryentry{microservices}{
    name={Microservices},
    description={Architekturmuster, bei dem eine Anwendung aus kleinen, unabhängig deploymbaren Services besteht}
}

\newglossaryentry{quality_gate}{
    name={Quality Gate},
    description={Definierte Qualitätskriterien in der Softwareentwicklung, die erfüllt sein müssen, bevor Code in die nächste Phase übergehen darf}
}

\newglossaryentry{kanban}{
    name={Kanban},
    description={Agile Methode zur Prozesssteuerung, die Arbeit visualisiert und den Arbeitsfluss durch Work-in-Progress-Limits optimiert}
}

\newglossaryentry{scrum}{
    name={Scrum},
    description={Agiles Framework für Projektmanagement, das in kurzen Iterationen (Sprints) arbeitet und auf selbstorganisierten Teams basiert}
}

\newglossaryentry{scope_creep}{
    name={Scope Creep},
    description={Unkontrollierte Erweiterung des Projektumfangs während der Entwicklung durch ungeplante Anforderungen, die zu Zeit- und Kostenüberschreitungen führen kann}
}

\newglossaryentry{skill}{
    name={Skill},
    plural={Skills},
    description={Fähigkeit oder Kompetenz eines Mitarbeiters, z.\,B. Programmiersprache, Technologie oder Methodik}
}

\newglossaryentry{three_tier}{
    name={Three-Tier-Architektur},
    description={Dreischichtige Softwarearchitektur bestehend aus Präsentationsschicht (Frontend), Anwendungsschicht (Backend/Geschäftslogik) und Datenhaltungsschicht (Datenbank)}
}

\newglossaryentry{wcag}{
    name={WCAG},
    description={Web Content Accessibility Guidelines -- Richtlinien für barrierefreie Webinhalte des W3C}
}